% Source: https://artofproblemsolving.com/community/c3h2530631_sum_of_all_real_solutions_x_polynomials
% Stable identity: aops-topic-2530631-q01
We need $x\ne0$, and multiplying by $x$ gives $x^4-2x^3+2x^2-x-15=0$. Notice that
\[x^4-2x^3+2x^2-x=\left(x^2-x\right)^2+\left(x^2-x\right).\]
So with $y=x^2-x$ the equation becomes $y^2+y-15=0$, i.e. $y=\dfrac{-1\pm\sqrt{61}}{2}$.

The equation $x^2-x=y$ has real solutions exactly when its discriminant $1+4y\ge0$, i.e. $y\ge-\dfrac14$.
\begin{itemize}
\item $y=\dfrac{-1+\sqrt{61}}{2}>-\dfrac14$: two distinct real solutions with sum $1$ (by Vieta for $x^2-x-y=0$).
\item $y=\dfrac{-1-\sqrt{61}}{2}<-\dfrac14$: no real solutions.
\end{itemize}
None of these solutions is $0$. So the sum of the real solutions is $\boxed{1}$.

Vieta's formulas for the quartic give the sum of \emph{all four} complex roots, $2$. The other two roots are the non-real solutions of $x^2-x=\dfrac{-1-\sqrt{61}}{2}$, and their sum, $1$, accounts for the difference.
